% LaTeX template for Artifact Evaluation V20240722
%
% Prepared by Grigori Fursin with contributions from Bruce Childers,
%   Michael Heroux, Michela Taufer and other colleagues.
%
% See examples of this Artifact Appendix in
%  * ASPLOS'24 "PyTorch 2: Faster Machine Learning Through Dynamic Python Bytecode Transformation and Graph Compilation": 
%      https://dl.acm.org/doi/10.1145/3620665.3640366
%  * SC'17 paper: https://dl.acm.org/citation.cfm?id=3126948
%  * CGO'17 paper: https://www.cl.cam.ac.uk/~sa614/papers/Software-Prefetching-CGO2017.pdf
%  * ACM ReQuEST-ASPLOS'18 paper: https://dl.acm.org/citation.cfm?doid=3229762.3229763
%
% (C)opyright 2014-2024 cTuning.org
%
% CC BY 4.0 license
%

\documentclass{sigplanconf}

\usepackage{hyperref}

\begin{document}

\special{papersize=8.5in,11in}

%%%%%%%%%%%%%%%%%%%%%%%%%%%%%%%%%%%%%%%%%%%%%%%%%%%%
% When adding this appendix to your paper, 
% please remove above part
%%%%%%%%%%%%%%%%%%%%%%%%%%%%%%%%%%%%%%%%%%%%%%%%%%%%

\appendix
\section{Artifact Appendix}

%%%%%%%%%%%%%%%%%%%%%%%%%%%%%%%%%%%%%%%%%%%%%%%%%%%%%%%%%%%%%%%%%%%%%
\subsection{Abstract}

This artifact contains the implementation of the translation and evaluation
pipeline used in the paper.  The artifact accepts Linux litmus tests and
RC11/C litmus tests written as \texttt{.litmus} files, translates them into a
common calculus litmus representation, and evaluates the translated test with
the executable memory-model checker in \texttt{calculus.py}.  The repository
also includes curated source corpora and regenerated translated corpora for the
Linux kernel tests, the Paul E. McKenney RCU tests, and the supported subset
of the RC11/C tests.  Running the workflow allows the reviewer to reproduce the
artifact's main outcome for an input test: whether the target execution is
\texttt{Allowed} or \texttt{Forbidden}.

\subsection{Artifact check-list (meta-information)}

{\em Obligatory. Use just a few informal keywords in all fields applicable to your artifacts
and remove the rest. This information is needed to find appropriate reviewers and gradually 
unify artifact meta information in Digital Libraries.}

{\small
\begin{itemize}
  \item {\bf Algorithm: } litmus-test translation, execution-graph exploration, memory-model checking
  \item {\bf Program: } Python 3 scripts, command-line workflow
  \item {\bf Compilation: } none
  \item {\bf Transformations: } Linux litmus to calculus litmus, RC11/C litmus to calculus litmus
  \item {\bf Binary: } none
  \item {\bf Model: } operational memory-model checker implemented in \texttt{calculus.py}
  \item {\bf Data set: } Linux kernel litmus tests, Paul E. McKenney RCU litmus tests, RC11/C litmus tests
  \item {\bf Run-time environment: } Python 3, local filesystem, command line
  \item {\bf Hardware: } standard x86-64 workstation or laptop
  \item {\bf Run-time state: } single-process command-line execution
  \item {\bf Execution: } local
  \item {\bf Metrics: } classification of executions as \texttt{Allowed} or \texttt{Forbidden}
  \item {\bf Output: } translated litmus file, final \texttt{Allowed}/\texttt{Forbidden} result
  \item {\bf Experiments: } single-test evaluation, corpus translation, corpus validation
  \item {\bf How much disk space required (approximately)?: } 100 MB
  \item {\bf How much time is needed to prepare workflow (approximately)?: } 5--10 minutes
  \item {\bf How much time is needed to complete experiments (approximately)?: } seconds for single tests; minutes for corpus checks
  \item {\bf Publicly available?: } yes, in the artifact branch of the repository
  \item {\bf Code licenses (if publicly available)?: } not specified in the repository
  \item {\bf Data licenses (if publicly available)?: } inherited from the included source corpora; not separately restated here
  \item {\bf Workflow automation framework used?: } none
  \item {\bf Archived (provide DOI)?: } not yet archived
\end{itemize}
}

%%%%%%%%%%%%%%%%%%%%%%%%%%%%%%%%%%%%%%%%%%%%%%%%%%%%%%%%%%%%%%%%%%%%%
\subsection{Description}

\subsubsection{How to access}

The artifact is the contents of this repository, using the
\texttt{artifact-eval} branch.  The main entry point is:

\begin{verbatim}
python3 scripts/run_artifact_pipeline.py path/to/test.litmus
\end{verbatim}

If \texttt{--kind} is omitted, the script prompts for the input language
(\texttt{linux} or \texttt{c}).  The repository also contains pre-generated
translations under \texttt{converted/}.

\subsubsection{Hardware dependencies}

No special hardware is required.  The artifact was prepared to run on a normal
CPU-based workstation or laptop.  GPU accelerators are not needed.

\subsubsection{Software dependencies}

The artifact requires Python 3.  During preparation it was run with
\texttt{Python 3.14.6}.  No external workflow framework is required.
The repository uses only the Python standard library plus the packages already
used by \texttt{calculus.py}.

\subsubsection{Data sets}

The repository includes the following litmus-test corpora:

\begin{itemize}
  \item \texttt{litmus/Kernel/}: Linux kernel litmus tests.
  \item \texttt{litmus/paulmckrcu/}: merged Paul E. McKenney RCU source corpus.
  \item \texttt{litmus/c/}: RC11/C litmus tests written as \texttt{.litmus} files.
\end{itemize}

Regenerated translated corpora are included under:
\texttt{converted/Kernel/}, \texttt{converted/paulmckrcu/}, and
\texttt{converted/c/}.

\subsubsection{Models}

The executable model is the memory-model checker implemented in
\texttt{calculus.py}.  The translator lowers source tests to a common event
language using \texttt{Read}, \texttt{Write}, \texttt{Fence}, and simple
control flow, after which \texttt{calculus.py} explores candidate relations and
reports whether the requested outcome is allowed.

%%%%%%%%%%%%%%%%%%%%%%%%%%%%%%%%%%%%%%%%%%%%%%%%%%%%%%%%%%%%%%%%%%%%%
\subsection{Installation}

No build step is required.  Clone or unpack the repository, enter the project
root, and verify that Python 3 is available:

\begin{verbatim}
python3 --version
\end{verbatim}

The artifact can then be run directly.  For example:

\begin{verbatim}
python3 scripts/run_artifact_pipeline.py \
  litmus/Kernel/MP+polocks.litmus --kind linux

python3 scripts/run_artifact_pipeline.py \
  litmus/c/a1+Racq+rel.litmus --kind c
\end{verbatim}

%%%%%%%%%%%%%%%%%%%%%%%%%%%%%%%%%%%%%%%%%%%%%%%%%%%%%%%%%%%%%%%%%%%%%
\subsection{Experiment workflow}

The standard workflow is:

\begin{enumerate}
  \item Choose a source litmus test from \texttt{litmus/Kernel/},
        \texttt{litmus/paulmckrcu/}, or \texttt{litmus/c/}.
  \item Run \texttt{scripts/run_artifact_pipeline.py} on the test.
  \item If desired, save the translated calculus litmus file with
        \texttt{--keep-translated}.
  \item Inspect the final output line, which reports
        \texttt{Allowed} or \texttt{Forbidden}.
\end{enumerate}

Representative commands:

\begin{verbatim}
python3 scripts/run_artifact_pipeline.py \
  litmus/Kernel/MP+unlocklockonceonce+fencermbonceonce.litmus --kind linux

python3 scripts/run_artifact_pipeline.py \
  litmus/paulmckrcu/C-LB-Lrw+R-A+R-A+R-OC.litmus --kind linux

python3 scripts/run_artifact_pipeline.py \
  litmus/c/arfna.litmus --kind c \
  --keep-translated /tmp/arfna.translated.litmus
\end{verbatim}

The artifact also supports evaluating already translated files directly:

\begin{verbatim}
python3 scripts/run_artifact_pipeline.py \
  converted/Kernel/MP+polocks.litmus --kind converted
\end{verbatim}

%%%%%%%%%%%%%%%%%%%%%%%%%%%%%%%%%%%%%%%%%%%%%%%%%%%%%%%%%%%%%%%%%%%%%
\subsection{Evaluation and expected results}

For a supported test, the expected final result is a single line ending with:

\begin{verbatim}
...: Allowed
\end{verbatim}

or

\begin{verbatim}
...: Forbidden
\end{verbatim}

Examples that should run successfully in the current artifact:

\begin{itemize}
  \item \texttt{litmus/Kernel/MP+unlocklockonceonce+fencermbonceonce.litmus}
  \item \texttt{litmus/paulmckrcu/C-LB-Lrw+R-A+R-A+R-OC.litmus}
  \item \texttt{litmus/c/arfna.litmus}
\end{itemize}

At the time this appendix was prepared, the refreshed translated corpora
contained:

\begin{itemize}
  \item 52 translated Linux kernel tests,
  \item 2768 translated Paul E. McKenney RCU tests,
  \item 106 translated RC11/C tests.
\end{itemize}

The remaining unsupported RC11/C tests fail clearly during translation rather
than silently producing malformed output.  The unsupported cases are currently
limited to atomic CAS/RMW patterns and arithmetic-over-read expressions such as
\texttt{linearisation.litmus} and \texttt{linearisation2.litmus}.

%%%%%%%%%%%%%%%%%%%%%%%%%%%%%%%%%%%%%%%%%%%%%%%%%%%%%%%%%%%%%%%%%%%%%
\subsection{Experiment customization}

Reviewers can supply any supported Linux or RC11/C \texttt{.litmus} input file.
The most useful customization points are:

\begin{itemize}
  \item choose a different source test in \texttt{litmus/},
  \item save the translated file with \texttt{--keep-translated},
  \item run on pre-generated files under \texttt{converted/} with
        \texttt{--kind converted}.
\end{itemize}

%%%%%%%%%%%%%%%%%%%%%%%%%%%%%%%%%%%%%%%%%%%%%%%%%%%%%%%%%%%%%%%%%%%%%
\subsection{Notes}

The artifact currently assumes that all user-supplied source inputs are
\texttt{.litmus} files.  When using source tests, the user must select the
source language as \texttt{linux} or \texttt{c}.  The repository intentionally
omits older manual \texttt{.c} translations and keeps the checked-in translated
corpora synchronized with the current translators.

%%%%%%%%%%%%%%%%%%%%%%%%%%%%%%%%%%%%%%%%%%%%%%%%%%%%%%%%%%%%%%%%%%%%%
\subsection{Methodology}

Submission, reviewing and badging methodology:

\begin{itemize}
  \item \url{https://www.acm.org/publications/policies/artifact-review-and-badging-current}
  \item \url{https://cTuning.org/ae}
\end{itemize}

%%%%%%%%%%%%%%%%%%%%%%%%%%%%%%%%%%%%%%%%%%%%%%%%%%%%
% When adding this appendix to your paper, 
% please remove below part
%%%%%%%%%%%%%%%%%%%%%%%%%%%%%%%%%%%%%%%%%%%%%%%%%%%%

\end{document}
